\documentclass[12pt,a4paper]{article}

\usepackage{cmap}
\usepackage[T2A]{fontenc}
\usepackage[utf8]{inputenc}
\usepackage[russian]{babel}
\usepackage{amsmath,amssymb}
\usepackage[a4paper,margin=2cm]{geometry}

\setlength{\parindent}{0pt}
\setlength{\parskip}{0.45em}
\sloppy

\begin{document}

\begin{center}
{\Large\bfseries Билет 10: вопросы для проверки знаний}
\end{center}

\noindent\fbox{%
  \begin{minipage}{\dimexpr\linewidth-2\fboxsep-2\fboxrule\relax}
  \normalfont
Числовые ряды. Критерий Коши сходимости ряда. Знакопостоянные ряды:
признак сравнения, признаки Даламбера и Коши, интегральный признак.
Знакопеременные ряды, сходимость и абсолютная сходимость. Признаки Лейбница,
Дирихле и Абеля. Независимость суммы абсолютно сходящегося ряда от порядка
слагаемых. Произведение абсолютно сходящихся рядов.
  \end{minipage}%
}

\section*{Вопросы на понимание}

\begin{enumerate}
\item Почему сходимость ряда определяется через частичные суммы, а не только через поведение его членов?
\item Приведите пример ряда, у которого общий член стремится к нулю, но сам ряд расходится. Что показывает этот пример?
\item Почему для ряда с неотрицательными членами достаточно проверить ограниченность частичных сумм сверху?
\item Что может измениться, а что не может измениться в сходимости ряда, если заменить конечное число первых членов?
\item Почему в признаке сравнения важно, чтобы сравниваемые ряды были знакопостоянными или неотрицательными начиная с некоторого места?
\item Почему признаки Даламбера и Коши не решают вопрос о сходимости в предельном случае, когда получается \(q=1\)?
\item Чем отличаются выводы признаков Даламбера и Коши при \(q<1\), \(q>1\) и \(q=1\)?
\item Какую роль играет монотонность функции в интегральном признаке для ряда, составленного из значений \(f(k)\)?
\item Чем условная сходимость отличается от абсолютной на примере знакопеременного ряда?
\item В признаке Дирихле почему одновременно нужны ограниченность частичных сумм одного множителя и стремление другого множителя к нулю?
\item Почему признак Лейбница можно рассматривать как частный случай признака Дирихле?
\item Почему для перестановки членов и произведения рядов существенна именно абсолютная сходимость?
\end{enumerate}

\section*{Источники для сверки}

\texttt{target/answer\_question\_10.tex};
\texttt{IvGE\_1.pdf}, стр. 264--280, 286--292;
\texttt{IvGE\_2.pdf}, стр. 286.

\end{document}
